\chapter{Progressions, Bohr neighbourhoods, and partition regularity}
\label{app:chapter}

This chapter expands Sections 3--5 and Appendices B--C of Ortega and
Prendiville, \emph{Extremal Sidon Sets are Fourier Uniform, with Applications
to Partition Regularity}, arXiv:2110.13447v1.  The arguments below use Haar
probability measure on $\mathbb T=\mathbb R/\mathbb Z$, the character
$e(t)=\exp(2\pi i t)$, and $[N]=\{1,\ldots,N\}\subset\mathbb Z$.
Every function on $[N]$ is extended by zero to $\mathbb Z$.  Throughout,
$N\geq1$, $m=|S|$, and
\[
 D=\sup_{\theta\in\mathbb T}
 \left|\widehat{1_S}(\theta)-\frac mN\widehat{1_{[N]}}(\theta)\right|,
 \qquad E=\left|\frac m{\sqrt N}-1\right|+N^{-1/6}.
\]
The Fourier-uniformity theorem supplies $D\leq C\sqrt N\sqrt E$ with an
absolute constant $C$.  The tensor-kernel proof in Appendix B of the paper
omits a dimension factor in a union bound.  We record the precise bound
it proves, then give a different smoothing argument that recovers the
absolute constants in both Bohr equidistribution corollaries.
The dimension-uniform bandwidth claimed in Lemma 4.1 cannot hold for
arbitrary maximum-metric Lipschitz functions: a counterexample is proved
below.  The exact Sidon energy and progression equidistribution are checked
in Lean.  Both Bohr equidistribution corollaries are checked with absolute
constants.  The corrected finite-bandwidth approximation, its
dimension-uniform counterexample, and the full transference and partition
regularity arguments are also checked in Lean.

\section{Arithmetic progressions}

\begin{lemma}[An $L^1$ estimate including the small cases]
\label{app:progression-l1}
\lean{Sidon.fourierL1_arithmeticProgression_le}\leanok
If $P\subset\mathbb Z$ is a finite arithmetic progression, then
\[
 \int_{\mathbb T}|\widehat{1_P}(\theta)|\,d\theta
 \leq 2+\log(|P|+1)\leq C_1\log(2+|P|)
\]
for an absolute constant $C_1$.
\end{lemma}
\begin{proof}
The empty progression has transform zero, and a singleton has transform of
modulus one.  These cases explain why the printed $\log|P|$ must be replaced
by $\log(2+|P|)$ in a statement valid for every finite progression.
For a nonsingleton progression write
$P=\{a+qj:0\leq j<L\}$, with $L\geq2$ and $q\ne0$.  Translation contributes
only the factor $e(a\theta)$, whose modulus is one.  Multiplication by the
nonzero integer $q$ preserves Haar probability measure on $\mathbb T$:
for continuous periodic $H$, split $[0,1]$ into $|q|$ equal intervals,
substitute $u=q\theta$ on each, and use periodicity.  Consequently its
$L^1$ norm equals that of $\sum_{j=0}^{L-1}e(j\theta)$.

The triangle inequality bounds this sum by $L$.  If
$\theta\ne0$ in $\mathbb T$, multiplying the sum by $1-e(\theta)$ gives
$1-e(L\theta)$, so its modulus is at most $2/|1-e(\theta)|$.
Represent $\theta$ by $t\in[-1/2,1/2]$.  Concavity of sine on
$[0,\pi/2]$ gives $|\sin(\pi t)|\geq2|t|$, and hence
$|1-e(\theta)|=2|\sin(\pi t)|\geq4\|\theta\|_{\mathbb T}$.
Thus the integrand is at most
$\min(L,(2\|\theta\|_{\mathbb T})^{-1})$ away from zero; at zero use $L$.
Integrating over $|t|\leq L^{-1}$ contributes at most $2$.
On $L^{-1}<|t|\leq1/2$ integrate $1/(2|t|)$, obtaining at most
$\log L$.  When $L=2$ the second region is empty.  Enlarging the numerical
constant covers every $L\geq0$ and proves the assertion.
\end{proof}

\begin{lemma}[Fourier discrepancy tested on a progression]
\label{app:progression-discrepancy}
\lean{Sidon.progression_discrepancy_le_fourier_norm}\leanok
\uses{app:progression-l1}
For any finite $S\subset[N]$ and finite arithmetic progression $P$,
\[
 \left||S\cap P|-\frac{|[N]\cap P|\,m}{N}\right|
 \leq D(2+\log(N+1))\leq C_1D\log(2+N).
\]
\end{lemma}
\begin{proof}
Put $Q=P\cap[N]$ and $f=1_S-(m/N)1_{[N]}$.  Intersecting an arithmetic
progression with an integer interval restricts its parameter $j$ to an
interval of consecutive integers; consequently $Q$ is again an arithmetic
progression, possibly empty or a singleton, and $|Q|\leq N$.
Expand both finite Fourier sums and integrate term by term.  The identity
$\int_{\mathbb T}e(k\theta)\,d\theta=1$ for $k=0$ and $0$ otherwise gives
\[
 \sum_x1_Q(x)f(x)
 =\int_{\mathbb T}\widehat{1_Q}(\theta)\widehat f(-\theta)\,d\theta.
\]
The left side is the discrepancy in the statement.  Take its absolute value,
move the norm inside the integral, bound $|\widehat f(-\theta)|$ by $D$,
and apply Lemma~\ref{app:progression-l1} and monotonicity of logarithm.
\end{proof}

\begin{theorem}[Paper Theorem 3.1]
\label{app:progression-equidistribution}
\lean{Sidon.sidon_progression_equidistribution,Sidon.sidon_progression_equidistribution_log}\leanok
\uses{app:progression-discrepancy}
For a Sidon set $S\subset[N]$, any finite arithmetic progression $P$, and
any $\varepsilon>0$,
\[
 |S\cap P|=\frac{|[N]\cap P|\,m}{N}
       +O_{\varepsilon}(N^{1/2+\varepsilon}E^{1/2}).
\]
\end{theorem}
\begin{proof}
Substitute $D\leq C\sqrt N\sqrt E$ in
Lemma~\ref{app:progression-discrepancy}.  To make the logarithm absorption
explicit, $\log N=\int_1^N t^{-1}\,dt\leq(N^{\varepsilon}-1)/\varepsilon$
for $N\geq1$, and $\log(2+N)\leq\log3+\log N$.
Thus $\log(2+N)\leq(\log3+1/\varepsilon)N^{\varepsilon}$.
No lower bound on $m$ is needed here; in particular $S=\varnothing$ is allowed.
\end{proof}

\begin{corollary}[Paper Corollaries 1.5 and 1.6]
\label{app:interval-residue-corollaries}
\lean{Sidon.sidon_interval_equidistribution,Sidon.sidon_residue_equidistribution}\leanok
\uses{app:progression-equidistribution}
Assume $m\geq\sqrt N/100$.  For any interval $I\subset[0,1]$, with any
choice of open or closed endpoints, and for $q\geq1$ and $a\in\mathbb Z$,
\[
 \frac1m\sum_{x\in S}1_I(x/N)
 =\operatorname{length}(I)+O_{\varepsilon}(N^{\varepsilon}E^{1/2}),
\]
\[
 \frac1m\sum_{x\in S}1_{a+q\mathbb Z}(x)
 =\frac1q+O_{\varepsilon}(N^{\varepsilon}E^{1/2}).
\]
The constants are independent of the endpoints, $a$, and $q$.
\end{corollary}
\begin{proof}
The integers $x\in[N]$ with $x/N\in I$ form an integer interval $P$.
Rounding each endpoint changes its expected length by at most one, so
$\big||P|/N-\operatorname{length}(I)\big|\leq2/N$; this also covers empty
and one-point intervals and the absence of $0$ from $[N]$.
Apply Theorem~\ref{app:progression-equidistribution} to $P$ and divide by
$m\geq\sqrt N/100$.  The additional error $2/N$ is absorbed because
$E^{1/2}\geq N^{-1/12}$ and $N^{-1}\leq N^{\varepsilon-1/12}$.

For residue classes, $P=[N]\cap(a+q\mathbb Z)$ is a finite progression.
Writing $N=qk+b$ with $0\leq b<q$ shows that its size is either $k$ or
$k+1$.  Therefore $\big||P|/N-1/q\big|\leq1/N$, including $q>N$.
Apply the same normalized discrepancy bound and absorb $1/N$ as above.
\end{proof}

\section{A quantitative trigonometric approximation}

\begin{definition}[Paper Definition B.1]
\label{app:torus-fourier}
For $d\geq1$, write $\mathbb T^d$ for the product torus with its Haar
probability measure, $\|\beta\|_{\max}=\max_j\|\beta_j\|_{\mathbb T}$,
and $m\cdot\beta=\sum_jm_j\beta_j$ for $m\in\mathbb Z^d$.
For integrable $F:\mathbb T^d\to\mathbb C$, define
\[
 \widehat F(m)=\int_{\mathbb T^d}F(\beta)e(-m\cdot\beta)\,d\beta.
\]
\end{definition}

\begin{lemma}[Tensor Fejer kernel]
\label{app:tensor-fejer}
\lean{Sidon.circleFejer,Sidon.integral_circleFejer,Sidon.circleFejerSpectrum_eq_max,Sidon.tensorFejer_eq_sum,Sidon.integral_tensorFejer_tail_le}\leanok
\uses{app:torus-fourier}
For an integer $M\geq1$, define
\[
 k_M(t)=\frac1M\left|\sum_{j=0}^{M-1} e(jt)\right|^2,
 \qquad \Lambda_M(\beta)=\prod_{j=1}^d k_M(\beta_j).
\]
Then $\Lambda_M\geq0$, $\int\Lambda_M=1$, and
\[
 \Lambda_M(\beta)=\sum_{m\in\mathbb Z^d}
       w_M(m)e(m\cdot\beta),\qquad
 w_M(m)=\prod_j(1-|m_j|/M)_+.
\]
Only $|m_j|<M$ contribute.  For $0<\eta$,
\[
 \int_{\|\beta\|_{\max}>\eta}\Lambda_M(\beta)\,d\beta
 \leq \frac d{M\eta^2}.
\]
\end{lemma}
\begin{proof}
Nonnegativity follows from the squared modulus formula.  Expanding that
square gives $M^{-1}\sum_{0\leq u,v<M}e((u-v)t)$.
For an integer $h$ with $|h|<M$, exactly $M-|h|$ pairs satisfy $u-v=h$;
for $|h|\geq M$ no pairs do.  This proves the one-dimensional expansion.
Taking the product proves the displayed expansion and its finite support.
Orthogonality leaves only $h=0$ in the integral of $k_M$, so its integral
is one.  Fubini then gives the corresponding product integral.

The geometric-series estimate from Lemma~\ref{app:progression-l1} implies
$k_M(t)\leq1/(M\|t\|_{\mathbb T}^2)$ when $t\ne0$.
On the set $\|\beta_j\|_{\mathbb T}>\eta$ this yields
\[
 \Lambda_M(\beta)\leq\frac1{M\eta^2}
                \prod_{i\ne j}k_M(\beta_i).
\]
Its integral over that set is at most $1/(M\eta^2)$, since the measure
in the remaining coordinate is at most one and all other factors have
integral one.  The set $\|\beta\|_{\max}>\eta$ is the union of these $d$
sets.  Nonnegativity and the union bound give the factor $d$.
If $\eta\geq1/2$, the set is empty and the same bound holds.
\end{proof}

\begin{lemma}[Paper Lemma 4.1, with explicit dimension dependence]
\label{app:trigonometric-approximation}
\lean{Sidon.exists_trigonometric_approximation,Sidon.ceil_fejer_bandwidth_le}\leanok
\uses{app:tensor-fejer}
Suppose $F:\mathbb T^d\to[0,1]$ is $K$-Lipschitz for the maximum metric,
$K\geq1$, and $\varepsilon>0$.  There is a trigonometric polynomial
$F_\varepsilon:\mathbb T^d\to[0,1]$ with
$\|F-F_\varepsilon\|_\infty\leq\varepsilon$, with frequencies
$|m_j|<M$, where for $0<\varepsilon\leq1$ one may take
\[
 M=\left\lceil27dK^2\varepsilon^{-3}\right\rceil.
\]
In particular $M\leq28dK^2\varepsilon^{-3}$ in this range.
For $\varepsilon\geq1$, the zero polynomial suffices.
\end{lemma}
\begin{proof}
A Lipschitz function is continuous and integrable on the compact torus.
For $M\geq1$ put
$F_M(\alpha)=\int F(\alpha-\beta)\Lambda_M(\beta)\,d\beta$.
Because the kernel is nonnegative with integral one, $0\leq F_M\leq1$.
Insert its finite character expansion, interchange a finite sum with the
integral, and substitute $u=\alpha-\beta$ using translation and inversion
invariance of Haar measure.  This gives
\[
 F_M(\alpha)=\sum_{|m_j|<M}w_M(m)\widehat F(m)e(m\cdot\alpha).
\]
The output is real because the convolution is real, or equivalently because
the coefficient at $-m$ is the conjugate of the coefficient at $m$.

Use mass one to express $F(\alpha)-F_M(\alpha)$ as the integral of
$(F(\alpha)-F(\alpha-\beta))\Lambda_M(\beta)$.
On $\|\beta\|_{\max}\leq\eta$, its absolute integral is at most $K\eta$.
On the complement, $|F(\alpha)-F(\alpha-\beta)|\leq1$ and
Lemma~\ref{app:tensor-fejer} bounds the integral by $d/(M\eta^2)$.
Choose $\eta=\varepsilon/(3K)>0$ and
$M=\lceil27dK^2/\varepsilon^3\rceil\geq1$.
The defining ceiling inequality gives
$27dK^2\leq M\varepsilon^3$, and therefore
\[
 K\eta=\varepsilon/3,\qquad
 \frac d{M\eta^2}=\frac{9dK^2}{M\varepsilon^2}
 \leq\varepsilon/3.
\]
Thus the error is at most $2\varepsilon/3\leq\varepsilon$.
Each coefficient has modulus at most one: the tensor weight lies in
$[0,1]$, and $|\widehat F(m)|\leq\int|F|\leq1$.
The ceiling adds at most one, which is at most
$dK^2\varepsilon^{-3}$ when $\varepsilon\leq1$.
For $\varepsilon\geq1$, $|F-0|\leq1$ proves the final assertion.
\end{proof}

\begin{proposition}[No dimension-independent finite bandwidth]
\label{app:no-uniform-bandwidth}
\lean{Sidon.BandwidthCounterexample.exists_real_interpolation_away_from_zero,Sidon.BandwidthCounterexample.exists_probability_weights_of_signed_moments,Sidon.BandwidthCounterexample.exists_continuous_dimension_obstruction,Sidon.BandwidthCounterexample.not_exists_uniform_dimension_bandwidth}\leanok
For every integer $M\geq0$, there are a positive dimension $d$ and a
$1$-Lipschitz function $F_d:\mathbb T^d\to[0,1]$, for the maximum metric,
such that no trigonometric polynomial with all coordinate frequencies
$|n_i|\leq M$ approximates $F_d$ uniformly within $1/16$.
Consequently the bandwidth bound in paper Lemma 4.1 cannot have an
absolute constant independent of $d$.
\end{proposition}
\begin{proof}
Fix $\rho=1/4$, put $L=2M+1$, and choose
$t_j=1/4+j/(4(L+1))$ for $0\leq j<L$.
These are distinct points in $[1/4,1/2)$, so their circle distances from
zero are at least $\rho$.  The numbers $z_j=e(t_j)$ are distinct because
the circle character $e$ is injective on $\mathbb R/\mathbb Z$.
Define the ordinary complex Lagrange polynomials
\[
 \ell_j(X)=\prod_{k\ne j}\frac{X-z_k}{z_j-z_k}.
\]
All denominators are nonzero, each degree is at most $L-1=2M$, and
$\ell_j(z_k)$ is one for $j=k$ and zero otherwise.
For any polynomial $Q$ of degree at most $2M$, the difference
$Q-\sum_j Q(z_j)\ell_j$ has degree at most $2M$ and vanishes at $L$
distinct points, hence is zero by the polynomial root bound.
Apply this identity to $Q=X^{n+M}$, where $-M\leq n\leq M$, and evaluate
at one.  With $\lambda_j=\ell_j(1)z_j^M$, the result is
\[
 \sum_j\lambda_j e(nt_j)=1\qquad(|n|\leq M).
\]
Use the same identity at $-n$ and take complex conjugates to obtain
$\sum_j\overline{\lambda_j}e(nt_j)=1$.
Averaging the two identities proves that the real numbers
$A_j=\operatorname{Re}\lambda_j$ satisfy
$\sum_j A_j e(nt_j)=1$ for every $|n|\leq M$.
Taking $n=0$ gives $\sum_jA_j=1$.
Put $A_j^+=\max(A_j,0)$, $A_j^-=\max(-A_j,0)$ and
$Z=\sum_jA_j^+=1+\sum_jA_j^-\geq1$.
The probability measures
\[
 \mu=Z^{-1}\left(\delta_0+\sum_jA_j^-\delta_{t_j}\right),
 \qquad \nu=Z^{-1}\sum_jA_j^+\delta_{t_j}
\]
have the same Fourier moments at each integer $|n|\leq M$:
$1+\sum_j A_j^-e(nt_j)=\sum_j A_j^+e(nt_j)$ before division by $Z$.
The measure $\mu$ assigns probability $a=1/Z>0$ to zero and its remaining
mass to points at distance at least $\rho$ from zero.
The whole support of $\nu$ has distance at least $\rho$ from zero.

For any positive integer $d$, put
\[
 F_d(\beta)=\max_{1\leq i\leq d}
                 (\rho-\|\beta_i\|_{\mathbb T})_+.
\]
The distance to zero, positive part, and maximum inequalities show that
$F_d$ is $1$-Lipschitz for the maximum metric.
Its range is $[0,\rho]\subset[0,1]$.
Under the product measure $\nu^{\otimes d}$ it vanishes identically.
Under $\mu^{\otimes d}$ it equals $\rho$ exactly when at least one
coordinate is zero and equals zero otherwise.  Independence gives
\[
 \int F_d\,d\mu^{\otimes d}-\int F_d\,d\nu^{\otimes d}
       =\rho\bigl(1-(1-a)^d\bigr).
\]
Choose $d$ so that $(1-a)^d<1/2$.  Such a finite positive $d$ exists
because $0\leq1-a<1$; if $a=1$, take $d=1$.
The last discrepancy is then strictly greater than $\rho/2=1/8$.

For any multi-index with $|n_i|\leq M$, the expectation of
$e(\sum_i n_i\beta_i)$ under either product measure factors into the
one-coordinate Fourier moments.  Those moments agree, so every polynomial
of the stated bandwidth has the same expectation under the two products.
If such a polynomial $P$ satisfied $\|F_d-P\|_\infty\leq1/16$, the
triangle inequality would bound the discrepancy of $F_d$ between these
probability measures by $2/16=1/8$, a contradiction.
The polynomial is not required to have range $[0,1]$, so the obstruction
also applies when that extra condition is imposed.

Finally, a dimension-independent bound $M\leq C K^2\varepsilon^{-3}$,
with $K=1$ and $\varepsilon=1/16$, would give one finite upper bandwidth
$\lceil4096C\rceil$ for every dimension.  Apply the construction at that
bandwidth to contradict the asserted approximation.  This proves the
failure of the absolute-constant interpretation of Lemma 4.1, while
leaving the dimension-dependent corrected lemma intact.
\end{proof}

\section{Smooth and regular Bohr neighbourhoods}

\begin{lemma}[A dimension-uniform smoothing alternative]
\label{app:compact-smoothing}
\lean{Sidon.exists_torus_smoothing,Sidon.norm_torusTest_le_compact_smoothing}\leanok
\uses{app:torus-fourier}
Let $d\geq1$, $K\geq1$, and $F:\mathbb T^d\to[0,1]$ be $K$-Lipschitz
for the maximum metric.  For every $0<\eta\leq1/2$ there is a continuous
$A_\eta:\mathbb T^d\to[0,1]$ such that
\[
 \|F-A_\eta\|_\infty\leq K\eta,\qquad
 A_\eta(\beta)=\sum_{u\in\mathbb Z^d}a_u e(u\cdot\beta),\qquad
 \sum_u|a_u|\leq(10/\eta)^d.
\]
The series converges absolutely and uniformly.  The coefficients are not
required to have finite support.
\end{lemma}
\begin{proof}
Define the continuous circle ramp
\[
 w(t)=\min(1,\max(0,((\eta/2)-\|t\|_{\mathbb T})/(\eta/4))),
 \qquad c=\int w,\qquad v=w/c.
\]
It takes values in $[0,1]$, equals one when $\|t\|\leq\eta/4$, and
vanishes when $\|t\|\geq\eta/2$.  A circle ball of radius
$r\leq1/2$ has Haar measure $2r$, by representing it in $[-1/2,1/2]$.
Thus $c\geq\eta/2>0$, and the normalization is legitimate.
Since $w^2\leq w$, we obtain
\[
 v\geq0,\quad\int v=1,\quad
 \int|v|^2=c^{-2}\int w^2\leq c^{-1}\leq2/\eta\leq10/\eta.
\]
Set $V(\beta)=\prod_jv(\beta_j)$.  Finite-product integration gives
$\int V=1$ and $\int|V|^2\leq(10/\eta)^d$.
If $\|\beta\|_{\max}>\eta/2$, some coordinate has norm greater than
$\eta/2$, and its factor is zero.  Thus $V$ has the asserted small support.

Put $V^*(\beta)=V(-\beta)$ and $g=V*V^*$, with
$(f*h)(x)=\int f(x-y)h(y)\,dy$.
These are continuous functions on a compact space, so all integrands,
including the product-space integrands below, are integrable.
Fubini and translation invariance give $g\geq0$ and $\int g=(\int V)^2=1$.
If both factors of the integrand defining $g(x)$ are nonzero, then
$\|x-y\|\leq\eta/2$ and $\|y\|\leq\eta/2$; the triangle inequality gives
$\|x\|\leq\eta$.  Consequently $g(x)=0$ outside that ball.

For $\chi_u(\beta)=e(u\cdot\beta)$ use the coefficient convention
$\widehat f(u)=\int\overline{\chi_u(\beta)}f(\beta)\,d\beta$.
In a convolution coefficient, apply Fubini and substitute $x=z+y$ in
the inner integral.  The character splits into its $z$ and $y$ factors,
proving $\widehat{f*h}(u)=\widehat f(u)\widehat h(u)$.
Reflection of a real function conjugates its coefficient.  It follows that
$\widehat g(u)=|\widehat V(u)|^2\geq0$.

The tensor characters are orthonormal by finite-product integration and
the one-dimensional character integral.  Their star algebra separates
points: if two torus points differ, a frequency supported at a differing
coordinate distinguishes them.  Stone--Weierstrass gives density in
continuous functions; density of continuous functions in $L^2$ then gives
a complete orthonormal system.  Parseval therefore proves
\[
 \sum_u|\widehat g(u)|=\sum_u|\widehat V(u)|^2
       =\int|V|^2\leq(10/\eta)^d.
\]
Summability is established before pointwise inversion: the Weierstrass
test gives a uniformly convergent continuous character series, while the
complete orthonormal expansion has the same $L^2$ limit.  The two limits
agree almost everywhere, and continuity and full support of Haar measure
make them agree everywhere.

Finally set $A_\eta=F*g$.  Positivity and unit mass put its values in
$[0,1]$.  On the support of $g$ the Lipschitz bound gives
$|F(x-y)-F(x)|\leq K\eta$; integrating proves the uniform approximation.
The coefficients are $a_u=\widehat F(u)\widehat g(u)$.
Probability normalization gives $|\widehat F(u)|\leq\int|F|\leq1$, so
$\sum_u|a_u|\leq\sum_u|\widehat g(u)|\leq(10/\eta)^d$.
The same uniform-convergence and $L^2$-uniqueness argument gives the
claimed character expansion of $A_\eta$.
\end{proof}

\begin{theorem}[Paper Corollary 1.7, with an absolute constant]
\label{app:smooth-bohr}
\lean{Sidon.sidon_smooth_equidistribution}\leanok
\uses{app:compact-smoothing}
Let $d\geq1$, $K\geq1$, $F:\mathbb T^d\to[0,1]$ be $K$-Lipschitz for
the maximum metric, and $\alpha\in\mathbb T^d$.  If $S\subset[N]$ is
Sidon and $m\geq\sqrt N/100$, then
\[
 \left|\frac1m\sum_{x\in S}F(x\alpha)
       -\frac1N\sum_{x=1}^NF(x\alpha)\right|
 \leq C_2 K^{2/3}E^{1/(8d)}.
\]
Here $C_2$ is an absolute constant, independent also of $d$; the checked
implementation permits $C_2=48000$.
\end{theorem}
\begin{proof}
Write $R=D/m$, which is defined because the lower bound on $m$ makes it
positive.  Apply Lemma~\ref{app:compact-smoothing} at scale $\eta$.
An absolutely convergent character series can be interchanged with each
finite normalized average.  Each frequency $u\in\mathbb Z^d$ is tested at
$u\cdot\alpha$, where its discrepancy is at most $R$.
The coefficient absolute sum therefore bounds the discrepancy of
$A_\eta$ by $R(10/\eta)^d$.
The two uniform approximation errors add at most $2K\eta$.
Thus the desired discrepancy is at most
\[
 2K\eta+R(10/\eta)^d.                                   \tag{*}
\]
If $R=0$, let $\eta$ tend to zero to get discrepancy zero.
Otherwise put $t=(10^dR/K)^{1/(d+1)}>0$.
When $t\leq1/2$, take $\eta=t$ in (*) to get $3Kt$.
When $t>1/2$, the trivial discrepancy bound one is at most $2Kt$ because
$K\geq1$.  In both cases the discrepancy is at most
\[
 30K^{d/(d+1)}R^{1/(d+1)}.
\]
Fourier uniformity gives $R\leq100C E^{1/2}$.
The factor $(100C)^{1/(d+1)}$ is bounded by an absolute constant, so the
last bound is at most $A K^{d/(d+1)}E^{1/(2d+2)}$ with absolute $A$.

It remains to compare exponents without introducing dimension-dependent
constants.  Put $B=K^{2/3}E^{1/(8d)}$.  If $B\geq1$, the trivial bound
one suffices.  If $B<1$, then $E<K^{-16d/3}\leq1$, and
\[
 \frac{K^{d/(d+1)}E^{1/(2d+2)}}{K^{2/3}E^{1/(8d)}}
 =K^{(d-2)/(3(d+1))}E^{(3d-1)/(8d(d+1))}
 \leq K^{-5d/(3(d+1))}\leq1.
\]
Thus the stronger estimate implies the claimed bound with the absolute
constant $C_2=\max(1,A)$ in every case.
\end{proof}

\begin{definition}[Regularity in a denominator-free form]
\label{app:regular-bohr}
\lean{Sidon.bohrSet,Sidon.IsRegularBohrSet,Sidon.isRegularBohrSet_iff_ratio}\leanok
For $d\geq1$, $\rho>0$, and $\alpha\in\mathbb T^d$, define
\[
 B(\alpha,\rho)=\{x\in[N]:\|x\alpha\|_{\max}\leq\rho\}.
\]
Call this finite Bohr set regular when, for every real $\kappa$ with
$|\kappa|\leq1/(100d)$,
\[
 \left||B(\alpha,(1+\kappa)\rho)|-|B(\alpha,\rho)|\right|
 \leq100d|\kappa|\,|B(\alpha,\rho)|.
\]
For nonempty $B$ this is exactly paper Definition 1.9 after multiplying by
$|B|$.  The denominator-free condition also handles an empty $B$ without
an undefined quotient.  All radii in this condition are positive.
\end{definition}

\begin{lemma}[Lipschitz cutoff functions]
\label{app:bohr-cutoffs}
\lean{Sidon.clippedNormRamp,Sidon.lipschitzWith_clippedNormRamp,Sidon.bohr_discrepancy_le_cutoff_errors}\leanok
Let $0<t<\rho$.  Define $u(\beta)=\|\beta\|_{\max}$ and
\[
 F_+(\beta)=\min(1,\max(0,(\rho+t-u(\beta))/t)),\quad
 F_-(\beta)=\min(1,\max(0,(\rho-u(\beta))/t)).
\]
Both functions take values in $[0,1]$ and have Lipschitz constant at most
$t^{-1}$.  At every integer $x\in[N]$,
\[
 1_{B(\alpha,\rho-t)}(x)\leq F_-(x\alpha)
 \leq1_{B(\alpha,\rho)}(x)
 \leq F_+(x\alpha)\leq1_{B(\alpha,\rho+t)}(x).
\]
\end{lemma}
\begin{proof}
The distance to zero on a metric space is $1$-Lipschitz, and taking a
maximum of finitely many coordinate distances preserves this Lipschitz
constant for the maximum metric.  Multiplying by $t^{-1}$ multiplies the
constant by $t^{-1}$, while clipping a real number to $[0,1]$ is
$1$-Lipschitz.  These observations prove the continuity estimates.
If $u\leq\rho-t$, then $F_-=1$; if $u>\rho$, then $F_-=0$.
If $u\leq\rho$, then $F_+=1$; if $u>\rho+t$, then $F_+=0$.
The inequalities at the boundary points follow directly from these
formulas, so the use of closed Bohr sets causes no endpoint exception.
\end{proof}

\begin{theorem}[Paper Corollary 1.11, with an absolute constant]
\label{app:regular-bohr-equidistribution}
\lean{Sidon.sidon_regular_bohr_equidistribution}\leanok
\uses{app:smooth-bohr,app:regular-bohr,app:bohr-cutoffs}
Under the hypotheses on $S$ above, if $B=B(\alpha,\rho)$ is regular, then
\[
 \left|\frac{|S\cap B|}{m}-\frac{|B|}{N}\right|
 \leq C_3 d\rho^{-1}E^{1/(14d)}.
\]
\end{theorem}
\begin{proof}
For $\rho\geq1/2$, every coordinate has distance at most $1/2$, so
$B=[N]$ and the discrepancy is zero.  Assume $0<\rho<1/2$.
For $E\geq1$, the assertion follows from the bound one after enlarging
$C_3$, since $d/\rho\geq2$.
For $0<E<1$, put $t=E^{3/(40d)}\in(0,1)$.
If $t>\rho/(100d)$, then
\[
 \frac d\rho E^{1/(14d)}
 >\frac1{100}E^{1/(14d)-3/(40d)}\geq\frac1{100},
\]
because $1/14-3/40=-1/280<0$.  The trivial bound one proves this case.

It remains to treat $t\leq\rho/(100d)$, which in particular gives
$t<\rho$ and $t^{-1}\geq1$.  Apply the preceding cutoff lemma.
Regularity at $\kappa=\pm t/\rho$ yields
\[
 \left|\frac{|B(\alpha,\rho\pm t)|}{N}-\frac{|B|}{N}\right|
 \leq100d\rho^{-1}t\frac{|B|}{N}\leq100d\rho^{-1}t.
\]
Apply Theorem~\ref{app:smooth-bohr} with $K=t^{-1}$ to each cutoff.
Using its upper inequality with $F_+$ bounds $|S\cap B|/m$ from above by
$|B|/N+100d\rho^{-1}t+C_2t^{-2/3}E^{1/(8d)}$.
Using its lower inequality with $F_-$ bounds the same quantity from below
by $|B|/N-100d\rho^{-1}t-C_2t^{-2/3}E^{1/(8d)}$.
The choice of $t$ gives
$t^{-2/3}E^{1/(8d)}=E^{3/(40d)}=t$.
Since $d/\rho\geq1$, the total error is at most
$(100+C_2)d\rho^{-1}E^{3/(40d)}$.
Finally $3/40>1/14$ and $E<1$, so this is at most the asserted bound.
\end{proof}

\section{The real Fourier norm and its dual}

\begin{definition}[Paper Lemma C.2: the norms]
\label{app:fourier-dual-norm}
\lean{Sidon.FourierSpace,Sidon.FourierSpace.transform,Sidon.FourierSpace.dualNorm}
\leanok
On the real vector space of functions $f:[N]\to\mathbb R$, set
\[
 \|f\|_F=\sup_{\alpha\in\mathbb T}|\widehat f(\alpha)|,
 \qquad
 \|\phi\|_F^*=\sup_{\|f\|_F\leq1}
             \left|\sum_{x=1}^Nf(x)\phi(x)\right|.
\]
\end{definition}

\begin{lemma}[Norm facts and the real modulation repair]
\label{app:fourier-dual-properties}
\lean{Sidon.FourierSpace.abs_pairing_le,Sidon.FourierSpace.dualNorm_mul_le,Sidon.FourierSpace.abs_apply_le_dualNorm,Sidon.FourierSpace.dualNorm_one,Sidon.FourierSpace.isCompact_fourier_closedBall,Sidon.FourierSpace.isCompact_dualNorm_ball}
\leanok
\uses{app:fourier-dual-norm}
The two expressions above are norms.  Their unit balls are compact, and
\[
 |\langle f,\phi\rangle|\leq\|f\|_F\|\phi\|_F^*,\quad
 \|\phi\psi\|_F^*\leq\|\phi\|_F^*\|\psi\|_F^*,\quad
 \|\phi\|_\infty\leq\|\phi\|_F^*,\quad \|1\|_F^*=1.
\]
All products here are pointwise products of real functions.
\end{lemma}
\begin{proof}
The supremum defining $\|f\|_F$ is finite because
$|\widehat f(\alpha)|\leq\sum_x|f(x)|$.  It is a maximum because the
transform is continuous on the compact torus.  Homogeneity and the triangle
inequality follow pointwise and then pass to the supremum.
Fourier inversion gives
$f(x)=\int\widehat f(\alpha)e(-x\alpha)\,d\alpha$, and hence
$|f(x)|\leq\|f\|_F$.  Thus norm zero forces $f=0$.
This estimate makes the unit ball bounded in the usual finite-dimensional
coordinate norm.  The Fourier norm is continuous (it is bounded above by
the coordinate $\ell^1$ norm and satisfies the reverse triangle inequality),
so its unit ball is closed and compact.

The pairing is continuous on that compact unit ball, so the dual supremum
is finite and attained.  If $f\ne0$, apply the definition to
$f/\|f\|_F$; if $f=0$, both sides of the pairing inequality are zero.
For each coordinate $x$, the delta function $1_{\{x\}}$ has Fourier norm
one.  Testing with it gives $|\phi(x)|\leq\|\phi\|_F^*$.
This also proves definiteness of the dual norm; its other norm axioms
follow from the supremum definition.

There is a small real-versus-complex issue in the printed proof of C.2:
the dual norm was defined using real test functions, so it cannot be applied
directly to $f(x)e(\alpha x)$.
For a correct real argument, fix $\alpha,\theta\in\mathbb T$ and put
$c(x)=\cos(2\pi(\alpha x+\theta))$.
The identity $c=(e(\theta)e(\alpha x)+e(-\theta)e(-\alpha x))/2$ gives
\[
 \widehat{fc}(\beta)=\tfrac12\left(
 e(\theta)\widehat f(\beta+\alpha)
 +e(-\theta)\widehat f(\beta-\alpha)\right),
 \qquad \|fc\|_F\leq\|f\|_F.
\]
Since $fc$ is real, the dual pairing inequality now applies to it.
Use $|z|=\sup_\theta\operatorname{Re}(e(\theta)z)$ with
$z=\widehat{f\phi}(\alpha)$ to obtain
$|\widehat{f\phi}(\alpha)|\leq\|f\|_F\|\phi\|_F^*$.
Taking the supremum in $\alpha$ gives
$\|f\phi\|_F\leq\|f\|_F\|\phi\|_F^*$.
Apply the pairing inequality to $f\phi$ and $\psi$, and then take the
supremum over real $f$ with $\|f\|_F\leq1$.  This proves
submultiplicativity with constant exactly one.
Finally $|\langle f,1\rangle|=|\widehat f(0)|\leq\|f\|_F$, while a delta
function realizes equality, so $\|1\|_F^*=1$.
\end{proof}

\begin{lemma}[Paper Lemma C.1: strict finite-dimensional separation]
\label{app:strict-separation}
\lean{Sidon.FourierSpace.strict_separation}
\leanok
If $K\subset\mathbb R^n$ is closed and convex and $v\notin K$, then there
is $\phi\in\mathbb R^n$ such that
$\langle v,\phi\rangle>\langle u,\phi\rangle$ for every $u\in K$.
\end{lemma}
\begin{proof}
If $K$ is empty, any $\phi$ works.  Otherwise choose $u_1\in K$.
The continuous distance to $v$ attains its minimum on
$K\cap\overline B(v,\|u_1-v\|+1)$, which is nonempty and compact.
Outside that ball the distance is larger than the distance to $u_1$, so a
minimizer $u_0$ there is also a global minimizer on $K$.
Set $\phi=v-u_0\ne0$.  For any $u\in K$ and $0<t\leq1$,
$u_0+t(u-u_0)\in K$; the minimality inequality expands to
\[
 -2t\langle\phi,u-u_0\rangle+t^2\|u-u_0\|^2\geq0.
\]
Divide by $t$ and let $t$ tend to zero from above.  This gives
$\langle\phi,u-u_0\rangle\leq0$.
Consequently $\langle v,\phi\rangle-\langle u,\phi\rangle
\geq\|v-u_0\|^2>0$, as required.
\end{proof}

\begin{lemma}[Polynomial approximation of a maximum, with a finite recipe]
\label{app:polynomial-max}
\lean{Sidon.exists_polynomial_positiveMax,Sidon.FourierSpace.exists_dualNorm_aeval_bound}
\leanok
For $r\geq1$, $L\geq1$ and $\tau>0$, a real polynomial $P$ in $r$
variables satisfies
\[
 |P(t)-\max(0,t_1,\ldots,t_r)|\leq\tau
 \quad\text{for }t\in[-L,L]^r.
\]
Its degree and coefficients can be chosen solely from $r,L,\tau$.
\end{lemma}
\begin{proof}
On $[0,1]^r$, let $H(z)=\max(0,2Lz_1-L,\ldots,2Lz_r-L)$.
It is $2L$-Lipschitz for the maximum metric.  Choose an integer
$k\geq\max(1,rL^2/\tau^2)$ and form its multivariate Bernstein polynomial
\[
 B_kH(z)=\sum_{a\in\{0,\ldots,k\}^r}H(a/k)
     \prod_{j=1}^r\binom{k}{a_j}z_j^{a_j}(1-z_j)^{k-a_j}.
\]
For fixed $z$, these nonnegative weights sum to one.  They are the
distribution of independent binomial random variables $Y_j$ with parameters
$(k,z_j)$.  The binomial identities for mean and variance give
$\mathbb E(Y_j/k-z_j)^2=z_j(1-z_j)/k\leq1/(4k)$.
These identities can equivalently be proved by differentiating
$(z+(1-z))^k$ once and twice and expanding finite sums.
By Cauchy--Schwarz and $\max_j b_j^2\leq\sum_jb_j^2$,
\[
 |B_kH(z)-H(z)|\leq2L\mathbb E\max_j|Y_j/k-z_j|
 \leq2L\sqrt{r/(4k)}\leq\tau.
\]
Substitute $z_j=(t_j+L)/(2L)$ to obtain $P$.
Its degree is at most $rk$, and the displayed finite formula fixes all its
coefficients using only $r,L,\tau$.  This supplies a constructive substitute
for the Stone--Weierstrass invocation in Appendix C.
\end{proof}

The checked implementation uses the existing real Stone--Weierstrass theorem
on the fixed compact cube $[-L,L]^r$.  It extracts an actual multivariate
polynomial from the range of the evaluation homomorphism.  Induction over
constants, addition, and multiplication by a variable then supplies a finite
dual-norm bound depending only on that polynomial and $L$.  Thus the checked
existence proof is uniform in $N$; the Bernstein formula above gives an
additional explicit informal witness.


\section{The colouring dense-model lemma}

\begin{theorem}[Paper Lemma 5.1, contrapositive form]
\label{app:dense-model}
\lean{Sidon.FourierSpace.colouring_dense_model,Sidon.colouring_dense_model_interval}
\leanok
\uses{app:fourier-dual-properties,app:strict-separation,app:polynomial-max}
For every integer $r\geq1$ and $0<\varepsilon\leq1$, there is
$\eta=\eta(\varepsilon,r)>0$ such that the following holds for every
$N\geq1$.  If $\nu,f_1,\ldots,f_r:[N]\to[0,\infty)$ satisfy
$\sum_if_i\leq\nu$ and
\[
 \|\nu-1\|_F\leq\eta N,
\]
there are $g_i:[N]\to[0,\infty)$ with $\sum_ig_i\leq1$ and
$\|f_i-g_i\|_F\leq\varepsilon N$ for every $i$.
\end{theorem}
\begin{proof}
We prove that failure of the conclusion forces
$\|\nu-1\|_F>\eta N$ for a fixed positive $\eta$.
If $\sum_x\nu(x)>2N$, evaluation at frequency zero gives
$\|\nu-1\|_F>N$, so assume $\sum_x\nu(x)\leq2N$.
Put $a=1+\varepsilon/2$ and $q=\varepsilon N/4$.
In the finite-dimensional real vector space of $r$ tuples of functions,
consider
\[
 K=\{(g_i+h_i)_i:g_i\geq0,\ \sum_ig_i\leq1,
                              \|h_i\|_F\leq q\}.
\]
The constraints on $g$ define a closed bounded set and are convex.
Each Fourier-norm ball for $h_i$ is compact and convex by
Lemma~\ref{app:fourier-dual-properties}.  Their product and its continuous
image under addition are compact, so $K$ is closed and convex.

The tuple $f/a$ lies outside $K$.  Indeed, if $f_i/a=g_i+h_i$, then
$f_i-g_i=(a-1)g_i+ah_i$.  Since $0\leq g_i\leq1$,
$\|g_i\|_F\leq N$, and $a\leq3/2$, the Fourier norm of this difference
is at most $\varepsilon N/2+3\varepsilon N/8
=7\varepsilon N/8<\varepsilon N$, contrary to failure of the conclusion.
Strict separation therefore provides a tuple $\phi$ satisfying
\[
 a^{-1}\sum_i\langle f_i,\phi_i\rangle
 >\sum_i\langle g_i,\phi_i\rangle+\sum_i\langle h_i,\phi_i\rangle
 \quad(g,h\text{ admissible}).
\]
Taking $g=h=0$ shows that $A=\sum_i\langle f_i,\phi_i\rangle>0$.
Multiply every $\phi_i$ by the positive scalar $aN/A$.  After this
normalization, $\sum_i\langle f_i,\phi_i\rangle=aN$, and the right side
of the last display is strictly less than $N$ for all admissible $g,h$.

Choose $g=0$, all but one $h_i=0$, and let that $h_i$ maximize its pairing
over its Fourier-norm ball.  Compactness and symmetry of the ball show that
this maximum is $q\|\phi_i\|_F^*$.  Thus
$\|\phi_i\|_F^*<N/q=4/\varepsilon$ and
$\|\phi_i\|_\infty\leq4/\varepsilon$.
Write $M(x)=\max(0,\phi_1(x),\ldots,\phi_r(x))$.
At each $x$, assign all mass one to the least index at which a positive
maximum is achieved, and assign zero to all coordinates when the maximum
is zero.  This defines an admissible $g$ with
$\sum_i\langle g_i,\phi_i\rangle=\sum_xM(x)<N$.
On the other hand, nonnegativity and $\sum_if_i\leq\nu$ imply
\[
 aN=\sum_i\langle f_i,\phi_i\rangle
 \leq\sum_x\nu(x)M(x).
\]
Subtracting the bound on $\sum_xM(x)$ gives
$\langle\nu-1,M\rangle>\varepsilon N/2$.

Apply Lemma~\ref{app:polynomial-max} with
$L=4/\varepsilon$ and $\tau=\varepsilon/100$ to obtain a fixed polynomial
$P(t)=\sum_b p_b t^b$.  Since
$\sum_x|\nu(x)-1|\leq\sum_x\nu(x)+N\leq3N$, replacing $M$ by
$P(\phi_1,\ldots,\phi_r)$ changes the pairing by at most
$3\varepsilon N/100$.  The resulting pairing is therefore greater than
$\varepsilon N/4$.
Repeated submultiplicativity, including $\|1\|_F^*=1$ for the constant
monomial, gives
\[
 \|P(\phi_1,\ldots,\phi_r)\|_F^*
 \leq B:=\sum_b|p_b|(4/\varepsilon)^{|b|}.
\]
Here the finite number $B$ depends only on $\varepsilon,r$, regardless of
$N,f,\nu$.  The dual pairing inequality gives
$\varepsilon N/4<\|\nu-1\|_F B$.
Taking $\eta=\min(1,\varepsilon/(4\max(1,B)))>0$ proves the claimed
contrapositive, and hence the theorem.  This choice also covers the
mass-greater-than-$2N$ case treated initially.
\end{proof}

\begin{lemma}[Completing a normalized dense model to a partition]
\label{app:dense-partition-model}
\uses{app:dense-model}
\lean{Sidon.complete_dense_model,Sidon.colouring_dense_partition_model_interval}
\leanok
If $r\geq1$, $\sum_i f_i=\nu$ on $[N]$, and
$\sum_{x\in[N]}\nu(x)=N$, the preceding theorem may additionally require
$\sum_i g_i(x)=1$ on $[N]$, with the same prescribed final Fourier error.
The functions $g_i$ are extended by zero off $[N]$.
\end{lemma}
\begin{proof}
Invoke the preceding theorem with error $\tau=\varepsilon/(r+1)$.
Write $D(x)=1-\sum_i g_i(x)\geq0$ on $[N]$ and set $D=0$ outside.
Evaluation at frequency zero gives
$\sum_{x\in[N]}(f_i(x)-g_i(x))\leq\tau N$ for each $i$.
Summing these inequalities and using the normalization gives
$\sum_xD(x)\leq r\tau N$.  Consequently
$\|D\|_F\leq\|D\|_1\leq r\tau N$.
Replace the first colour by $g_1+D$ and leave the others unchanged.
All colours remain nonnegative and their pointwise sum becomes one on $[N]$;
each is therefore at most one.  The first Fourier error is at most
$(r+1)\tau N=\varepsilon N$, and every other error is at most $\tau N$.
\end{proof}

\section{Counting solutions and transferring a colouring}

\begin{definition}[Weighted linear solution count]
\label{app:linear-count}
\lean{Sidon.linearSolutions,Sidon.weightedSolutionCount,Sidon.integral_prod_weightedFourier}\leanok
Let $c=(c_1,\ldots,c_s)$ be nonzero integers.  For functions on $[N]$ put
\[
 T_c(h_1,\ldots,h_s)=
 \sum_{\substack{x\in[N]^s\\\sum_ic_ix_i=0}}\prod_{i=1}^s h_i(x_i),
 \qquad T_c(h)=T_c(h,\ldots,h).
\]
The variables may repeat.  The conclusions below do not assert pairwise
distinctness, which is not a condition in paper Theorem 1.15.
\end{definition}

\begin{lemma}[The quantitative colouring theorem]
\label{app:fgr-dependency}
\lean{Sidon.QuantitativeRado.exists_many_monochromatic_linearSolutions}\leanok
\uses{qr:many-monochromatic}
Suppose $c_i\ne0$, $s\geq1$, $r\geq1$, and some nonempty
$I\subset\{1,\ldots,s\}$ has $\sum_{i\in I}c_i=0$.
There are $N_0=N_0(c,r)$ and $\delta_0=\delta_0(c,r)>0$ such that every
$r$-colouring of $[N]$, for $N\geq N_0$, has a class $C_j$ with
$T_c(1_{C_j})\geq\delta_0N^{s-1}$.
\end{lemma}
\begin{proof}
This is Theorem~\ref{qr:many-monochromatic}, whose full proof in the next
chapter starts from the Hales--Jewett theorem already checked in mathlib.
It is the single-equation instance of Frankl--Graham--R\"odl,
\emph{Quantitative theorems for regular systems of equations}, Journal of
Combinatorial Theory A 47 (1988), Theorem 1, cited as reference [11] in the
paper.  The zero-sum subcollection is Rado's columns condition for a
one-row matrix.  Rado's existence theorem alone would not give the needed
$N^{s-1}$ solutions; the full-dimensional embedding and multiplicity count
in the next chapter supply that additional argument.
\end{proof}

\begin{lemma}[Paper Lemma 5.2: the weighted reduction]
\label{app:weighted-interval-count}
\lean{Sidon.QuantitativeRado.exists_many_weighted_linearSolutions}\leanok
\uses{app:fgr-dependency,app:linear-count}
For the same $c,r$, let $g_i:[N]\to[0,\infty)$ satisfy $\sum_ig_i=1$.
For $N\geq N_0$, some $j$ satisfies
\[
 T_c(g_j)\geq\delta N^{s-1},\qquad \delta=\delta_0/r^s>0.
\]
\end{lemma}
\begin{proof}
At each $x$, at least one of the $r$ values $g_i(x)$ is at least $1/r$:
otherwise their sum would be strictly less than one.  Choose the least such
index and let $C_i$ be the set assigned to $i$.  These sets partition $[N]$,
and $g_i\geq r^{-1}1_{C_i}$ pointwise.  Each term in a solution count is
a product of nonnegative numbers, so
$T_c(g_i)\geq r^{-s}T_c(1_{C_i})$.
Apply the quantitative colouring theorem to the partition and
substitute the definition of $\delta$.
\end{proof}

\begin{lemma}[The exact Sidon energy]
\label{app:sidon-energy}
\lean{IsSidonSet.addEnergy_eq,IsSidonSet.addEnergy_le}\leanok
\uses{def:sidon_set}
For a finite Sidon set $S\subset\mathbb Z$ of size $m$, its additive energy
is exactly
\[
 \#\{(a,b,c,d)\in S^4:a-b=c-d\}=2m^2-m\leq2m^2.
\]
\end{lemma}
\begin{proof}
For the zero difference, there are $m$ choices for each ordered pair and
therefore $m^2$ quadruples.  For each nonzero difference, the Sidon property
allows at most one ordered pair.  There are exactly $m(m-1)$ nonzero
ordered pairs, with distinct differences, so these contribute $m(m-1)$.
This proves the exact energy formula, including $m=0$.
The checked declarations use mathlib's natural-valued
$\operatorname{Finset.addEnergy}(S,S)$; the tuple equation above is the
equivalent difference formulation after rearranging the summands.
\end{proof}

\begin{lemma}[Fourth and higher Fourier moments]
\label{app:sidon-fourth-moment}
\uses{app:sidon-energy,app:linear-count}
If $S\subset[N]$ is Sidon, $m>0$, $p\in\mathbb N$ with $p\geq4$, and
$0\leq f\leq(N/m)1_S$, then
\[
 \int_{\mathbb T}|\widehat f(\theta)|^p\,d\theta
 \leq2N^{p}m^{-2}.
\]
If $p\in\mathbb N$ with $p\geq2$ and $0\leq g\leq1_{[N]}$, then the corresponding integral
is at most $N^{p-1}$.
\end{lemma}
\begin{proof}
The triangle inequality and the bound on $f$ give
$\|\widehat f\|_\infty\leq\sum_xf(x)\leq N$.
Expand $|\widehat f|^4=\widehat f\,\overline{\widehat f}
\widehat f\,\overline{\widehat f}$ and integrate term by term.
Orthogonality restricts the sum to $a-b+c-d=0$.
Every weight is nonnegative and at most $(N/m)^4$, so the fourth moment
is at most $(N/m)^4(2m^2-m)\leq2N^4/m^2$.
Pointwise $|\widehat f|^p\leq N^{p-4}|\widehat f|^4$, and integration
proves the displayed estimate.
For $g$, similarly $\|\widehat g\|_\infty\leq N$; Parseval and $g^2\leq g$
give $\int|\widehat g|^2=\sum_xg(x)^2\leq N$.
Multiply by $N^{p-2}$ to obtain $N^{p-1}$.
\end{proof}

\begin{lemma}[The five-variable counting inequality]
\label{app:counting-transfer}
\lean{Sidon.sidon_counting_transfer}\leanok
\uses{app:sidon-fourth-moment}
Suppose $s\geq5$, all $c_i\ne0$, $m\geq\sqrt N/100$, $S\subset[N]$ is
Sidon, $0\leq f\leq(N/m)1_S$, and $0\leq g\leq1_{[N]}$.  Then
\[
 |T_c(f)-T_c(g)|\leq40000s\,N^{s-2}\|f-g\|_F.
\]
\end{lemma}
\begin{proof}
For finite complex-valued functions $h_i$, expand their transforms and use
orthogonality to obtain
\[
 T_c(h_1,\ldots,h_s)
   =\int_{\mathbb T}\prod_i\widehat h_i(c_i\theta)\,d\theta.
\]
The map $\theta\mapsto c_i\theta$ preserves Haar measure because
$c_i\ne0$, as in Lemma~\ref{app:progression-l1}.
Consequently composing a Fourier transform with that map preserves every
finite $L^p$ norm.

Put $F_i(\theta)=\widehat f(c_i\theta)$,
$G_i(\theta)=\widehat g(c_i\theta)$, and
$b_i(\theta)=\max(|F_i(\theta)|,|G_i(\theta)|)$.
Expanding one factor at a time and using the triangle inequality proves
\[
 \left|\prod_iF_i-\prod_iG_i\right|
 \leq\sum_i|F_i-G_i|\prod_{j\ne i}b_j
 \leq\|f-g\|_F\sum_i\prod_{j\ne i}b_j.
\]
For completeness, inserting a factor replaces the difference by
$(F_i-G_i)\prod F+G_i(\prod F-\prod G)$; induction gives exactly the
displayed sum over omitted indices.
Set $p=s-1\geq4$.  The preceding moment lemma gives
\[
 \int|\widehat f|^{s-1}\leq2N^{s-1}/m^2\leq20000N^{s-2},
 \qquad \int|\widehat g|^{s-1}\leq N^{s-2}.
\]
Pointwise $b_j^p\leq|F_j|^p+|G_j|^p$, hence
$\int b_j^p\leq20001N^{s-2}$.
Apply the arithmetic--geometric mean inequality to the $p$ nonnegative
numbers $b_j^p$ with $j\ne i$:
\[
 \prod_{j\ne i}b_j\leq\frac1p\sum_{j\ne i}b_j^p.
\]
Integrating gives $\int\prod_{j\ne i}b_j\leq20001N^{s-2}$.
There are $s$ omitted-index terms; their sum is bounded by
$20001sN^{s-2}\|f-g\|_F$, which is at most the stated bound.
This identifies the exact use of $s\geq5$: four controlled moments must
remain after the distinguished factor is removed.
\end{proof}

\begin{theorem}[Paper Theorem 1.15, with explicit quantifier order]
\label{app:partition-regularity}
\lean{Sidon.sidon_partition_regularity_relative,Sidon.sidon_partition_regularity}
\leanok
\uses{app:dense-partition-model,app:weighted-interval-count,app:counting-transfer}
Let $s\geq5$, $r\geq1$, and $c_i\in\mathbb Z\setminus\{0\}$, and
suppose some nonempty subcollection of the coefficients sums to zero.
There are constants $N_*=N_*(c,r)\geq1$, $\gamma=\gamma(c,r)>0$, and
$\kappa=\kappa(c,r)>0$ such that for every $N\geq N_*$ and Sidon
$S\subset[N]$ with
\[
 \left|\frac{|S|}{\sqrt N}-1\right|<\gamma,
\]
every $r$-colouring $S=C_1\sqcup\cdots\sqcup C_r$ has a class satisfying
\[
 T_c(1_{C_j})\geq\kappa\frac{|S|^s}{N}.
\]
\end{theorem}
\begin{proof}
Let $N_0,\delta>0$ be the constants from the weighted interval lemma.
Set
\[
 \varepsilon=\min(1,\delta/(80000s))>0,
\]
and choose the positive
parameter $\eta$ from Lemma~\ref{app:dense-partition-model} for this
accuracy and these $r$ colours.  All three choices depend only on $c,r$.
Choose
\[
 \gamma=\min(1/2,\eta^2/4096)>0.
\]
Since $N^{-1/6}\to0$, choose $N_*\geq\max(1,N_0)$ such that
$N^{-1/6}\leq\eta^2/4096$ for every $N\geq N_*$.
These choices are made before $N,S$, and the colouring.

Write $m=|S|$ and $E=|m/\sqrt N-1|+N^{-1/6}$.
The deviation hypothesis and $\gamma\leq1/2$ give
$m\geq\sqrt N/2>0$, hence also $m\geq\sqrt N/100$.
Define $\nu=(N/m)1_S$ and $f_i=(N/m)1_{C_i}$, extended by zero
outside $[N]$.  These functions are nonnegative,
$\sum_i f_i=\nu$, and $\sum_x\nu(x)=N$.
The baseline Fourier theorem with constant $16$ gives
\[
 \|\nu-1_{[N]}\|_F
 =\frac Nm\left\|1_S-\frac mN1_{[N]}\right\|_F
 \leq32N\sqrt E.
\]
Here $E\leq2\eta^2/4096\leq\eta^2/1024$, so
$32\sqrt E\leq\eta$.
The complete dense-model lemma therefore supplies nonnegative $g_i$,
zero outside $[N]$, with $\sum_i g_i=1$ on $[N]$, $g_i\leq1$, and
$\|f_i-g_i\|_F\leq\varepsilon N$ for every $i$.

The weighted interval lemma supplies a colour $j$ for which
$T_c(g_j)\geq\delta N^{s-1}$.
The counting-transfer theorem applies because
$0\leq f_j\leq N/m$, its support is contained in $S$, and
$0\leq g_j\leq1$.
It gives
\[
 |T_c(f_j)-T_c(g_j)|
 \leq40000sN^{s-2}\|f_j-g_j\|_F
 \leq40000s\varepsilon N^{s-1}
 \leq\tfrac\delta2N^{s-1}.
\]
Consequently $T_c(f_j)\geq(\delta/2)N^{s-1}$.
By homogeneity of the $s$ factors,
$T_c(f_j)=(N/m)^s T_c(1_{C_j})$.
Multiplication by $(m/N)^s>0$ gives
\[
 T_c(1_{C_j})\geq\frac\delta2\frac{m^s}{N}.
\]
Thus $\kappa=\delta/2>0$ has the required property.
The checked cardinality form counts the same finite tuples: a product of
indicator values is one exactly when every coordinate belongs to $C_j$,
and zero otherwise.  Repeated coordinates remain permitted.
\end{proof}

\section{The introduction's obstruction examples}

\begin{proposition}[Paper Proposition 1.13]
\label{app:odd-obstruction}
\uses{construction:sharp-obstructions}
\lean{Sidon.eventually_exists_sharp_odd_obstruction}\leanok
There are Sidon sets $S_N\subset[N]$ with
$|S_N|\geq(1/\sqrt2+o(1))\sqrt N$ and with no solution in $S_N$ to
$x_1+x_2+x_3+x_4=x_5$.
\end{proposition}
\begin{proof}
Fix $\delta>0$. For every sufficiently large $N$, apply
Lemma~\ref{construction:relative-primes} with $x=\sqrt{N/2}$ and
$\varepsilon=\delta\sqrt2$. It supplies a prime $p$ with
\[
 (1/\sqrt2-\delta)\sqrt N\le p\le\sqrt{N/2}.
\]
Theorem~\ref{construction:bose-chowla} supplies a Sidon set
$S_0\subset[p^2-1]$ of cardinality $p$.
Set $S_N=\{2x+1:x\in S_0\}$. Its elements lie between $3$ and
$2(p^2-1)+1\le2p^2\le N$.
Equality of two sums in $S_N$ cancels two added ones and then the
nonzero factor two, so the original Sidon property proves that $S_N$
is Sidon. Injectivity gives $|S_N|=p$.
All its elements are odd. A sum of four odd integers is even, while
an element of $S_N$ is odd, ruling out the equation with repeated
variables allowed. The displayed lower bound holds for every $\delta>0$
once $N$ exceeds its associated threshold; this is the stated asymptotic.
\end{proof}

\begin{proposition}[The construction in paper Remark 1.14]
\label{app:nonzero-sum-obstruction}
\uses{construction:sharp-obstructions}
\lean{Sidon.exists_eventually_sharp_residue_obstruction}\leanok
If $c_1,\ldots,c_s\in\mathbb Z$ have $\sigma=\sum_ic_i\ne0$, there are
Sidon sets $S_N\subset[N]$ of size at least
$(p^{-1/2}+o(1))\sqrt N$, for some fixed prime $p$, containing no solution
to $\sum_ic_ix_i=0$.
\end{proposition}
\begin{proof}
Choose a prime $p>|\sigma|$ using the unboundedness of primes.
It cannot divide the nonzero integer $\sigma$: divisibility would force
$p\le|\sigma|$.
Fix $\delta>0$. The relative-prime theorem applied at $\sqrt{N/p}$
with parameter $\delta\sqrt p$ supplies, for all sufficiently large $N$,
a construction prime $\ell$ with
\[
 (p^{-1/2}-\delta)\sqrt N\le\ell\le\sqrt{N/p}.
\]
Take $S_0\subset[\ell^2-1]$ of size $\ell$ from Bose--Chowla and put
$S_N=pS_0+1$. Its largest possible element is
$p(\ell^2-1)+1\le p\ell^2\le N$; its smallest element is positive.
Cancellation proves the Sidon property and injectivity preserves size.
Every element is one modulo $p$, hence
$\sum_i c_ix_i\equiv\sigma\not\equiv0\pmod p$.
The lower cardinality bound holds for every $\delta>0$ eventually,
which proves the asymptotic assertion.
\end{proof}
